% Source: sources/returns/S6_04_MAPS_MANIFOLDS/sections/manifolds.tex:3–29
\paragraph{Theorem D.1.1 (Refuted: the proposed connected sum in Lead 5).}
Put $d=|12\ell_0-4\ell_1-3\ell_2|$ for an admissible member $X_p$.
The universal cover $\widetilde X_p$ is a compact complex threefold of
Euler characteristic $2d$. Its Betti numbers necessarily satisfy
\[
 b_3(\widetilde X_p)=2+2b_2(\widetilde X_p)-2d,
 \qquad b_2(\widetilde X_p)\geq d-1.
\]
Consequently, for $d>1$ its second integral homology cannot vanish,
and it cannot be diffeomorphic to $\#^{d-1}(S^3\times S^3)$.
The conditional clause of Lead 5 has an impossible premise for $d>1$;
it is not a classification of the covers.
\begin{proof}
The covering has $d$ sheets, is unramified, and pulls back the complex
atlas, so is compact, complex, and oriented. Multiplicativity of Euler
characteristic for a finite covering follows by lifting a finite CW
structure: every cell has exactly $d$ lifts. Thus
$\chi(\widetilde X_p)=d\chi(X_p)=2d$.
Simple connectivity gives $b_1=b_5=0$, and oriented Poincare duality gives
$b_4=b_2$, $b_0=b_6=1$. Substitution gives the displayed identity. The
inequality follows from $b_3\geq0$. Finally a connected sum of $r$
copies of $S^3\times S^3$ has Euler characteristic $2-2r$; for $r=d-1$
this equals $4-2d$, which equals $2d$ only when $d=1$.
For a concrete admissible example take $(\ell_0,\ell_1,\ell_2)=(0,2,-1)$.
Then $d=5$, the cover has Euler characteristic $10$, whereas
$\#^4(S^3\times S^3)$ has Euler characteristic $-6$.
\end{proof}
